% slides/08_source.tex


\begin{frame}{Solving the EEG Inverse Problem}
  \begin{columns}[T]
    \begin{column}{0.50\textwidth}
      \textbf{The inverse problem is ill-posed}
      \begin{itemize}
        \item Infinitely many source configurations can produce the same scalp topography
        \item Regularisation and anatomical priors are essential
      \end{itemize}
      \vspace{0.4em}
      \textbf{LORETA family} \cite{PascualMarqui1994}
      \begin{itemize}
        \item \textbf{LORETA}: Laplacian smoothness prior
        \item \textbf{sLORETA}: zero localisation error for single dipoles
        \item \textbf{eLORETA}: exact, standardised; gold standard for resting-state
      \end{itemize}
    \end{column}
    \begin{column}{0.47\textwidth}
      \textbf{Beamforming}
      \begin{itemize}
        \item \textbf{LCMV}: scalar/vector; optimal for ERPs
        \item \textbf{DICS}: frequency-domain; ideal for oscillatory sources \cite{Gross2001}
        \item Requires accurate covariance — sensitive to artifact residuals
      \end{itemize}
      \vspace{0.3em}
      \textbf{MRI-constrained methods}
      \begin{itemize}
        \item Individual head models (BEM, FEM)
        \item Cortical surface constraint (FreeSurfer parcellations)
        \item Combined EEG-fMRI asymptotic approaches
      \end{itemize}
    \end{column}
  \end{columns}
\end{frame}
